\documentclass[a4paper,12pt]{book} %
\usepackage[margin=2.5cm,a4paper]{geometry} %lmargin=2.5cm,rmargin=4cm,tmargin=3cm,bmargin=3cm

\usepackage{amsmath,amssymb,mathtools,lineno,graphicx,textcomp,booktabs,url,setspace,xcolor,soul,eurosym}
% \usepackage{multirow,listings,setspace,gnuplottex,latexsym,keyval,ifthen,moreverb,lscape,forest}
\usepackage{multirow,listings,setspace,latexsym,keyval,ifthen,moreverb,lscape,forest}

\graphicspath{{fig/}}

\usepackage[figuresright]{rotating}

\usepackage[]{microtype} %activate={true,nocompatibility},final,kerning=true,spacing=true,factor=1100,stretch=10,shrink=10

\usepackage{tikz,pgfplots}
\usetikzlibrary{calc,shapes,arrows,intersections,shadows}
\usepgflibrary{arrows.new} 
\usetikzlibrary{arrows.new} % LATEX and plain TEX when using TikZ
\tikzset{>=latex}

\newif\ifpienumberinlegend
\pgfkeys{/number in legend/.code=
    \expandafter\let\expandafter\ifpienumberinlegend
    \csname if#1\endcsname
    \ifpienumberinlegend
    \let\legendbeforenumber\beforenumber
    \let\legendafternumber\afternumber
    \def\beforenumber##1\afternumber{}%
    \fi,
    /number in legend/.default=true
}

\usepackage{chemfig} 

\usepackage[]{threeparttable}

\usepackage[indention=0.5cm,labelsep=colon,font={sf,small},labelfont={bf,sf}]{caption}
\usepackage[indention=0.5cm,font={sf,small},labelfont={bf,sf}]{subcaption}

\definecolor{light-gray}{gray}{0.90} %not always visible on a dell screen!!!


\usepackage{fancyhdr}  
\pagestyle{fancy} 
\renewcommand{\chaptermark}[1]{\markboth{#1}{}}
\fancypagestyle{frontmatter}{%    
	\fancyhf{} 	
	\fancyfoot[RO,LE]{\textsf{Page \thepage}} 
	\fancypagestyle{plain}{
		\renewcommand{\headrulewidth}{0pt}
		\fancyhead{}
		\fancyfoot[RO,LE]{\textsf{Page \thepage}}
	} 
}
\fancypagestyle{mainmatter}{%    
	\fancyhf{} 	
	\fancyhead[RO,LE]{\nouppercase{\small \textsf{\leftmark}}}
	\fancyfoot[RO,LE]{\textsf{Page \thepage}} 
}
\setlength{\headheight}{15pt}


\newcounter{letternum}
\newcounter{lettersum}
\setcounter{lettersum}{13}
\newlength{\thumbtopmargin}
\setlength{\thumbtopmargin}{1cm}
\newlength{\thumbbottommargin}
\setlength{\thumbbottommargin}{3cm}
\newlength{\thumbheight}
\pgfmathsetlength{\thumbheight}{%
(\paperheight-\thumbtopmargin-\thumbbottommargin)/\value{lettersum}}

\newlength{\thumbwidth}
\setlength{\thumbwidth}{1.5cm}

\tikzset{
   thumb/.style={
%   draw=black,
  fill=light-gray,
   text=black,
   minimum height=\thumbheight, %\thumbheight,
   text width=\thumbwidth,
	outer sep=0pt,%   outer sep=10pt,
   font=\sffamily\Huge,
}
}
\newcommand{\oddthumb}[1]{%
      \begin{tikzpicture}[remember picture, overlay]
        \node [thumb,text centered,anchor=north east,] at ($%
            (current page.north east)-%
            (0,\thumbtopmargin+\value{letternum}*\thumbheight)%
        $) {#1};
   \end{tikzpicture}
}
\newcommand{\eventhumb}[1]{%
      \begin{tikzpicture}[remember picture, overlay]
        \node [thumb,text centered,anchor=north west,] at ($%
            (current page.north west)-%
            (0,\thumbtopmargin+\value{letternum}*\thumbheight)%
        $) {#1};
   \end{tikzpicture}
}
% create a new command to set a new lettergroup
\newcommand{\lettergroup}[1]{%
\fancyhead[LO]{\oddthumb{#1}}%
\fancyhead[RE]{\eventhumb{#1}}%
\fancypagestyle{chapterstart}{%
    \renewcommand{\headrulewidth}{0pt}
    \renewcommand{\footrulewidth}{0pt}
    \fancyhf{}
    \chead{\oddthumb{#1}}% chapters start only on odd pages
    \fancyfoot[RO,LE]{\textsf{Page \thepage}} %\fancyfoot[RO,LE]{\textsf{Page \thepage}}
  }
    \thispagestyle{chapterstart}
\stepcounter{letternum}%
}


\usepackage[numbib,nottoc]{tocbibind}	%numbib

%\usepackage[final]{pdfpages}

\usepackage[]{natbib} %numbers,sort&compress

\usepackage[Sonny]{fncychap} %Sonny

\usepackage[colorlinks,linkcolor=black,urlcolor=black,citecolor=black]{hyperref} %load hyperref after fncychap
\hypersetup{%
    pdftitle = {A title},
    pdfsubject = {A subject},
    pdfkeywords = {import keywords},
    pdfauthor = {Me of course},
    pdfcreator = {\LaTeX\ with package \flqq hyperref\frqq},
}
\usepackage[]{cleveref} % load cleveref after hyperref

\setlength{\parindent}{0cm} 
\renewcommand{\contentsname}{Table of Contents}
\renewcommand{\listfigurename}{List of Figures}
\renewcommand{\listtablename}{List of Tables}
\renewcommand{\appendixname}{}

\makeatletter
\newenvironment{chapquote}[2][2em]
  {\setlength{\@tempdima}{#1}%
   \def\chapquote@author{#2}%
   \parshape 1 \@tempdima \dimexpr\textwidth-2\@tempdima\relax%
   \itshape}
  {\par\normalfont\hfill--\ \chapquote@author\hspace*{\@tempdima}\par\bigskip}
\makeatother

\setlength{\parindent}{0cm} 
\pgfplotsset{compat=1.12}

%%%%%%%%%%%%%%%% BEGIN DOCUMENT %%%%%%%%%%%%%%%%%%%%%%%
\begin{document} 
\onehalfspace

\frontmatter
\pagestyle{frontmatter}
\lstset{language=Perl}
%%%%%%%%%%%%%%%%  BEGIN TITLEPAGE  %%%%%%%%%%%%%%%%%%
\begin{titlepage}
\centering
\includegraphics[width=\textwidth]{DI-eps-converted-to.pdf}

\vspace{1cm}

\includegraphics[width=2.25cm]{ITGlogo_new} 
\vspace{2.25cm}

{\Huge \textsf{Proteomic analysis of \textit{Taenia} spp. excretion/secretion proteins: the search for \textit{Taenia solium}-specific diagnostic antigens}}

\vspace{2.25cm}

{\LARGE \textsf{Bjorn Victor}}

\vspace{2.25cm}

{\textsf{Proefschrift voorgedragen tot het behalen van de graad van \\Doctor in de Diergeneeskundige Wetenschappen, 2014}}

\vfill

{\textsf{Promotoren: \\
Prof. Dr. Pierre Dorny\\Dr. Sarah Gabri\"el\\ Dr. Magnus Palmblad}}

\vfill

{\textsf{Vakgroep Virologie, Parasitologie en Immunologie\\
Faculteit Diergeneeskunde, Universiteit Gent \\
Salisburylaan 133, B-9820 Merelbeke}}
\end{titlepage}

%%%%%%%%%%%%%%%%  END TITLEPAGE  %%%%%%%%%%%%%%%%%%

\newpage{\thispagestyle{empty}\cleardoublepage}

\thispagestyle{empty}
\phantom{text}

\vfill

\begin{chapquote}{Douglas Adams, \textit{The Long Dark Tea-Time of the Soul} (1988).}
``I may not have gone where I intended to go, but I think I have ended up where I needed to be.''\\
\end{chapquote}

\vfill

\newpage{\thispagestyle{empty}\cleardoublepage}
\chapter{Acknowledgements}
text

\newpage{\thispagestyle{empty}\cleardoublepage}
%%%%%%%%%%%%%%%%  BEGIN LISTS   %%%%%%%%%%%%%%%%%%%%%%%
\newpage{\thispagestyle{empty}\cleardoublepage}
\tableofcontents 

\newpage{\thispagestyle{empty}\cleardoublepage}
\listoffigures
  
\newpage{\thispagestyle{empty}\cleardoublepage}
\listoftables

\newpage{\thispagestyle{empty}\cleardoublepage}
%%%%%%%%%%%%%%%%  END LISTS   %%%%%%%%%%%%%%%%%%%%%%%
\mainmatter
%\linenumbers %this will destroy the red frame around the protein info!!!


%%%%%%%%%%%%%%%%%%%%%%% General introduction %%%%%%%%%%%%%%%%


\newpage{\thispagestyle{empty}\cleardoublepage}
\pagestyle{mainmatter}
\chapter{General Introduction} \label{chap:generalintro}
\lettergroup{\thechapter}
text

%%%%%%%%%%%%%%%%%%%%%%% T. solium literature %%%%%%%%%%%%%%%%


%
\newpage{\thispagestyle{empty}\cleardoublepage}
\chapter{\textit{Taenia solium}: An Introduction} \label{chap:tsolliterature}
\lettergroup{\thechapter} 

\section{\textit{Taenia solium} taeniasis and cysticercosis}
\textit{Taenia solium} is a tapeworm belonging to the Taeniidae family (\Cref{fig:taxo}).
\begin{sidewaysfigure}
\resizebox{1\linewidth}{!}{
\centering
\begin{forest}
[Kingdom [Phylum [Class [Order [Family [Genus\textit{\phantom{j}}]]]]]]
\end{forest}
\quad
\begin{forest}
[Metazoa
 [Platyhelminthes 
                      [Cestoda [Cyclophyllidea [Taeniidae [\textbf{\textit{Taenia}}][\textit{Echinococcus}]][Hymenolepididae[\textit{Hymenolepis}]]] ] 
                      [Trematoda [Strigeidida[Schistosomatidae[\textit{Schistosoma}]]] [Opisthorchiida[Opisthorchiidae[\textit{Clonorchis}]]]] 
                   ] 
 [Nematoda
                      [Enoplea[Trichocephalida[Trichinellidae[\textit{Trichinella}]][Trichuridae[\textit{Trichuris}]]] ]
                      [Chromadorea [Ascaridida[Ascarididae[\textit{Ascaris}]]][Rhabditida[Strongylidae[\textit{Strongylus}]][Rhabditidae[\textit{Caenorhabditis}]]] ]
                  ]
]
\end{forest}
}
\caption{text}
\label{fig:taxo}
\end{sidewaysfigure}
The adult stage of the worm \citep{Linne:1758uq} develops in the upper third section of the small intestine of humans, %which is the duodenum-jejunum
the only final host, where it causes an infection called taeniasis (\Cref{fig:lifecycle}). 
\begin{figure}[!htbp]
\centering
{\sffamily
\begin{tikzpicture} 
\tikzstyle{every node}=[anchor=center]
\tikzstyle{column 3}=[text width=4cm,text height=2.3cm,text centered]
\tikzstyle{column 2}=[text width=5cm,text centered]
\tikzstyle{column 1}=[text width=4cm,text height=1.6cm,text centered]
\matrix [row sep=-0.30cm,column sep=1cm] %draw 
{
\node(v0)[text height=0.6cm]{\large\underline{Intermediate host}}; 	&	&\node(v0)[text height=0.6cm]{\large\underline{Final host}}; \\ 
 &  \node(v2)[]{\scalebox{0.15}{ %0.6
\begin{tikzpicture}[line width=1mm] %use thick for all figures!!!
\tikzstyle{every node}=[anchor=west]
%\draw[help lines] (0,0) grid (15,10);
\draw [blue] (1,7) to [out=10,in=170] (14,7);
\draw [blue] (1,6) to [out=10,in=170] (14,6);
\draw [blue] (1,5) to [out=10,in=185] (5.9,5.65);
\draw [blue] (9.1,5.65) to [out=-5,in=170] (14,5);
\draw [blue] (1,4) to [out=10,in=185] (5,4.6);
\draw [blue] (10,4.6) to [out=-5,in=170] (14,4);
\draw [blue] (1,3) to [out=10,in=185] (5.5,3.6);
\draw [blue] (9.5,3.6) to [out=-5,in=170] (14,3);
\draw [blue] (1,2) to [out=10,in=170] (14,2);
\draw [blue] (1,1) to [out=10,in=170] (14,1);
\draw [blue,] (7.5,4.5) circle [radius=0.3];
\draw [blue] (7.2,4.55) to [out=75,in=-45] (7,5.96) ;
\draw [blue] (7.8,4.45) to [out=75,in=-65] (8,5.96) ;
\draw [blue,thick] (7.4,4.4) circle [radius=0.05];
\draw [blue,thick] (7.6,4.6) circle [radius=0.05];
\draw [blue,thick] (7.6,4.4) circle [radius=0.05];
\draw [blue,thick] (7.4,4.6) circle [radius=0.05];
\draw [blue] (7.5,4.5) ellipse (2.5cm and 1.5cm);
\end{tikzpicture}}\\Consumption of raw or undercooked infected pork};  		&\\ 

\node(v6){\scalebox{0.15}{ %0.6
\begin{tikzpicture}[line width=1mm] %use thick for all figures!!!
\tikzstyle{every node}=[anchor=west]
%\draw[help lines] (0,0) grid (15,10);

\draw [blue] (13,0) -- (13,5) to [out=90,in=0] (10,8) -- (5,8);% to [out=0,in=105] (8.75,0) -- (6.25,0);
\draw[blue] (13,0) -- (12.5,0.5) -- (12,0) -- (12,3) to [out=90,in=90] (11,3) -- (11,0) -- (10.5,0.5) -- (10,0) -- (10,3) to [out=210,in=330] (5,3) -- (5,0) -- (4.5,0.5) -- (4,0) -- (4,3) to [out=90,in=90] (3,3) -- (3,0) -- (2.5,0.5) -- (2,0) -- (2,5);

\draw [blue] (3,7) circle [radius=2.25]; %kop
\draw [blue] (3,6) circle [radius=0.75]; %neus
\draw [blue] (2.7,6) circle [radius=0.1];
\draw [blue] (3.30,6) circle [radius=0.1];
\draw [blue] (2,7.5) circle [radius=0.15]; %ogen
%\draw [blue,fill=blue] (2,7.45) circle [radius=0.07];
\draw [blue] (4,7.5) circle [radius=0.15]; %ogen
%\draw [blue,fill=blue] (4,7.45) circle [radius=0.07];

\draw [blue] (14,9) -- (12.11,7.11); %staart
\draw [blue] (12.7,8.35) circle [radius=0.45];

\draw [blue] (1,8) to [out=135,in=270] (0.5,9.5) to [out=0,in=135] (2,9) ;
\draw [blue] (4,9) to [out=45,in=180] (5.5,9.5) to [out=-90,in=45] (5,8) ;

\draw [blue] (17.5,5.5) ellipse (2cm and 1cm);
\draw [blue,] (17.5,5.5) circle [radius=0.3];
\draw [blue] (17.2,5.55) to [out=75,in=-45] (17,6.45) ;
\draw [blue] (17.8,5.45) to [out=75,in=-65] (18,6.45) ;
\draw [blue,thick] (17.4,5.4) circle [radius=0.05];
\draw [blue,thick] (17.6,5.6) circle [radius=0.05];
\draw [blue,thick] (17.6,5.4) circle [radius=0.05];
\draw [blue,thick] (17.4,5.6) circle [radius=0.05];

\end{tikzpicture}}\\Porcine cysticercosis};& &
\node(v4)[]{\scalebox{0.18}{ %0.6
\begin{tikzpicture}[line width=1mm] %use thick for all figures!!!
\tikzstyle{every node}=[anchor=west]
%\draw[help lines] (0,0) grid (15,15);

\draw [blue] (7.65,5) -- (7.65,1) to [out=270,in=180] (8.525,0) to [out=0,in=270] (9.4,1) -- (9.4,8) -- (9.5,8) -- (9.5,5) to [out=270,in=180] (10,4.5) to [out=0,in=270] (10.5,5) -- (10.5,9) to [out=90,in=0] (9.5,10) -- (5.5,10) to [out=180,in=90] (4.5,9) -- (4.5,5) to [out=270,in=180] (5,4.5) to [out=0,in=270] (5.5,5) -- (5.5,8) -- (5.6,8) -- (5.6,1) to [out=270,in=180] (6.475,0) to [out=0,in=270] (7.35,1) -- (7.35,5) -- (7.65,5);

\draw [blue] (7.5,11.8) circle [radius=1.3];

\draw [blue] (2,8) circle [radius=0.4];
\draw [blue,thick] (2.15,8.15) circle [radius=0.05];
\draw [blue,thick] (2.15,7.85) circle [radius=0.05];
\draw [blue,thick] (1.85,8.15) circle [radius=0.05];
\draw [blue,thick] (1.85,7.85) circle [radius=0.05];

\draw [blue] (2.2,7.65) to [out=-75,in=75] (2.2,4.6) ;
\draw [blue] (1.8,7.65) to [out=-75,in=75] (1.8,4.6) ;
\draw [blue] (2.2,4.6) to [out=-105,in=105] (2.2,2) ;
\draw [blue] (1.8,4.6) to [out=-105,in=105] (1.8,2) ;

\draw [blue] (1.95,7) -- (2.35,7) ;
\draw [blue] (2.05,6) -- (2.45,6) ;
\draw [blue] (1.9,5) -- (2.3,5) ;
\draw [blue] (1.65,4) -- (2.05,4) ;
\draw [blue] (1.6,3) -- (2,3) ;

\end{tikzpicture}}\\Human taeniasis}; 	\\ 

	& \node(v8)[]{\scalebox{0.1}{ %0.6
\begin{tikzpicture}[line width=1mm] %use thick for all figures!!!
\tikzstyle{every node}=[anchor=west]
%\draw[help lines] (0,0) grid (15,10);

\draw [blue] (7.5,5) circle [radius=2.5];
\draw [blue] (7.5,5) circle [radius=1.8];
\draw [blue] (7.35,5) to [out=90,in=-45] (7,6) ;
\draw [blue] (7.65,5) to [out=90,in=-45] (7.4,6) ;
\draw [blue] (7,4.5) to [out=135,in=0] (6.25,5) ;
\draw [blue] (7.3,4.6) to [out=135,in=0] (6.5,5.2) ;
\draw [blue] (7.7,4.6) to [out=45,in=180] (8.25,5.2) ;
\draw [blue] (8.,4.5) to [out=45,in=180] (8.5,5) ;
\end{tikzpicture}}\\Eggs ingested by pigs\\ (coprophagy or accidental)}; 			&\\ 
\node(v12){\scalebox{0.18}{ %0.6
\begin{tikzpicture}[line width=1mm] %use thick for all figures!!!
\tikzstyle{every node}=[anchor=west]
%\draw[help lines] (0,0) grid (15,15);

\draw [blue] (7.65,5) -- (7.65,1) to [out=270,in=180] (8.525,0) to [out=0,in=270] (9.4,1) -- (9.4,8) -- (9.5,8) -- (9.5,5) to [out=270,in=180] (10,4.5) to [out=0,in=270] (10.5,5) -- (10.5,9) to [out=90,in=0] (9.5,10) -- (5.5,10) to [out=180,in=90] (4.5,9) -- (4.5,5) to [out=270,in=180] (5,4.5) to [out=0,in=270] (5.5,5) -- (5.5,8) -- (5.6,8) -- (5.6,1) to [out=270,in=180] (6.475,0) to [out=0,in=270] (7.35,1) -- (7.35,5) -- (7.65,5);
\draw [blue] (7.5,11.8) circle [radius=1.3];

\draw [blue] (15.5,5.5) ellipse (1.7cm and 0.9cm);
\draw [blue,] (15.5,5.5) circle [radius=0.3];
\draw [blue] (15.2,5.55) to [out=75,in=-45] (15,6.35) ;
\draw [blue] (15.8,5.45) to [out=75,in=-65] (16,6.35) ;
\draw [blue,thick] (15.4,5.4) circle [radius=0.05];
\draw [blue,thick] (15.6,5.6) circle [radius=0.05];
\draw [blue,thick] (15.6,5.4) circle [radius=0.05];
\draw [blue,thick] (15.4,5.6) circle [radius=0.05];

\end{tikzpicture}}\\Human (neuro) cysticercosis};	&  \node(v11){\scalebox{0.1}{ %0.6
\begin{tikzpicture}[line width=1mm] %use thick for all figures!!!
\tikzstyle{every node}=[anchor=west]
%\draw[help lines] (0,0) grid (15,10);
\draw [blue] (7.5,5) circle [radius=2.5];
\draw [blue] (7.5,5) circle [radius=1.8];
\draw [blue] (7.35,5) to [out=90,in=-45] (7,6) ;
\draw [blue] (7.65,5) to [out=90,in=-45] (7.4,6) ;
\draw [blue] (7,4.5) to [out=135,in=0] (6.25,5) ;
\draw [blue] (7.3,4.6) to [out=135,in=0] (6.5,5.2) ;
\draw [blue] (7.7,4.6) to [out=45,in=180] (8.25,5.2) ;
\draw [blue] (8.,4.5) to [out=45,in=180] (8.5,5) ;
\end{tikzpicture}}\\Eggs (accidentally) ingested by humans}; 		&\\
}; 
%\draw [->] (v2.west) -- (v4.north) node [above,text centered,midway] {replacement};
\draw [-latex new, arrow head=2mm] [blue,thick] (v2.east) -- (v4.north);
%\draw [arrows=->] [blue,thick] (v1.south) -- (v4.north);
\draw [-latex new, arrow head=2mm] [blue,thick] (v4.south) -- (v8.east);
\draw [-latex new, arrow head=2mm] [blue,thick] (v4.south) -- (v11.east) node [right,blue,midway,text width=2.8cm,text height=1.2cm] {\textcolor{black}{Proglottids/eggs \\shed with stools}};
%\draw [arrows=->] [blue,thick] (v9.north) -- (v6.south);
\draw [-latex new, arrow head=2mm] [blue,thick] (v8.west) -- (v6.south);
\draw [-latex new, arrow head=2mm] [blue,thick] (v11.west) -- (v12.east);
\draw [-latex new, arrow head=2mm] [blue,thick] (v6.north) -- (v2.west);
\end{tikzpicture}}
\caption{text}
\label{fig:lifecycle}
\end{figure}



%%%%%%%%%%%%%%%%%% The Hitchhiker's Guide to Proteomics %%%%%%%%%%%%%%%%



\newpage{\thispagestyle{empty}\cleardoublepage}
\chapter[The Hitchhiker's Guide to Proteomics]{The Hitchhiker's Guide to Proteomics\protect\footnote{Title adapted from the novel ``The Hitchhiker's Guide to the Galaxy'', written by Douglas Adams and published by Pan Books in 1979.}}\label{chap:proteomics}
\lettergroup{\thechapter}

\begin{figure}[!htbp]
\centering
{\sffamily
\begin{tikzpicture} 
\tikzstyle{every node}=[anchor=center]
\tikzstyle{column 1}=[text width=2.5cm,text centered]
\tikzstyle{column 2}=[text width=3cm,text centered]
\tikzstyle{column 3}=[text width=3cm,text centered]
\tikzstyle{column 4}=[text width=3cm,text centered]
\matrix [row sep=0.25cm,column sep=1cm] %draw,
{
\node(v1)[]{genomics}; 	& \node(v2)[]{transcriptomics};& \node(v3)[]{PROTEOMICS};		&\node(v4)[]{metabolomics};\\
\node(v5)[]{what can happen}; 		& \node(v6)[]{what will \\happen};		& 	\node(v7)[]{what is happening};								&\node(v8)[]{what has happened};\\
}; 
\draw [-latex new, arrow head=3mm] [thick] (v1.east) -- (v2.west);
\draw [-latex new, arrow head=3mm] [thick] (v2.east) -- (v3.west);
\draw [-latex new, arrow head=3mm] [thick] (v3.east) -- (v4.west);
\draw [arrows=-] [thick,dotted] (v1.south) -- (v5.north);
\draw [arrows=-] [thick,dotted] (v2.south) -- (v6.north);
\draw [arrows=-] [thick,dotted] (v3.south) -- (v7.north);
\draw [arrows=-] [thick,dotted] (v4.south) -- (v8.north);
\end{tikzpicture}}
\caption{text}
\label{fig:proteomics}
\end{figure}


\begin{figure}[!htbp]
\centering
\subcaptionbox{IEF. \label{fig:IEF}}[.35\linewidth]{\scalebox{0.65}{
\begin{tikzpicture}[very thick] %use thick for all figures!!!
\tikzstyle{every node}=[anchor=west]
%\draw[help lines] (0,0) grid (15,4);

\draw [black] (-3,2.25) -- (4,2.25) -- (4,1.75) -- (-3,1.75) -- (-3,2.25);
\draw [black] (-2,2.25) -- (-2,1.75);
\draw [black] (-1,2.25) -- (-1,1.75);
\draw [black] (-1.35,2.25) -- (-1.35,1.75);
\draw [black] (0,2.25) -- (0,1.75);
\draw [black] (0.15,2.25) -- (0.15,1.75);
\draw [black] (0.3,2.25) -- (0.3,1.75);
\draw [black] (1,2.25) -- (1,1.75);
\draw [black] (2,2.25) -- (2,1.75);
\draw [black] (2.2,2.25) -- (2.2,1.75);
\draw [black] (3,2.25) -- (3,1.75);
\draw [black] (3.3,2.25) -- (3.3,1.75);
\draw [black] (3.4,2.25) -- (3.4,1.75);
\draw [black] (3.6,2.25) -- (3.6,1.75);
\node[anchor=south,black,rotate=0] at (-2.5,2.4) {\Large \sffamily Low pH};
\node[anchor=south,black,rotate=0] at (-3,0.6) {\Large \sffamily (+)};
\draw [-latex new, arrow head=3mm,black] (-2.3,1) -- (0,1);
\draw [-latex new, arrow head=3mm,black] (3.3,1) -- (1,1);
\node[anchor=south,black,rotate=0] at (0.5,0.6) {\Large \sffamily pI};
\node[anchor=south,black,rotate=0] at (3.5,2.4) {\Large \sffamily High pH};
\node[anchor=south,black,rotate=0] at (4,0.6) {\Large \sffamily ($-$)};
\node[black,rotate=0] at (2.5,0.2) {\Large \sffamily \phantom{aa}};

\end{tikzpicture}}}\hfill
\subcaptionbox{1D-PAGE. \label{fig:1-Delectr}}[.3\linewidth]{\scalebox{0.65}{
\begin{tikzpicture}[very thick] %use thick for all figures!!!
\tikzstyle{every node}=[anchor=west]
%\draw[help lines] (0,0) grid (15,4);
\draw [black] (7,0) -- (7,4) -- (7.5,4) -- (7.5,3.5) -- (8,3.5 ) -- (8,4) -- (8.5,4) -- (8.5,3.5) -- (9,3.5) -- (9,4) -- (9.5,4) -- (9.5,3.5) -- (10,3.5) -- (10,4) -- (10.5,4) -- (10.5,3.5) -- (11,3.5) -- (11,4) -- (11.5,4) -- (11.5,3.5) -- (12,3.5) -- (12,4) -- (12.5,4) -- (12.5,0) -- (7,0);

\draw [thick,black] (7.5,2.2) -- (8,2.2);
\draw [thick,black] (7.5,1.5) -- (8,1.5);
\draw [thick,black] (7.5,0.75) -- (8,0.75);
\draw [thick,black] (8.5,3) -- (9,3);
\draw [thick,black] (8.5,2) -- (9,2);
\draw [thick,black] (8.5,1) -- (9,1);
\draw [thick,black] (9.5,2.5) -- (10,2.5);
\draw [thick,black] (9.5,2.8) -- (10,2.8);
\draw [thick,black] (9.5,1.5) -- (10,1.5);
\draw [thick,black] (10.5,2.7) -- (11,2.7);
\draw [thick,black] (10.5,1.8) -- (11,1.8);
\draw [thick,black] (10.5,0.5) -- (11,0.5);
\draw [thick,black] (11.5,3) -- (12,3);
\draw [thick,black] (11.5,1.5) -- (12,1.5);
\draw [thick,black] (11.5,0.75) -- (12,0.75);

\node[anchor=south,black,rotate=0] at (13.25,3) {\Large \sffamily ($-$)};
\node[anchor=south,black,rotate=0] at (13.25,0) {\Large \sffamily (+)};
\draw [-latex new, arrow head=3mm,black] (13.25,3) -- (13.25,1) node [above,black,midway,rotate=90] {\Large \sffamily size};

%\draw [-latex new, arrow head=3mm] [thick] (7,0) -- (8,0) node [above,black,midway] {pH$\uparrow$};

\end{tikzpicture}}}\hfill
\subcaptionbox{2D-PAGE. \label{fig:2-Delectr}}[.3\linewidth]{\scalebox{0.65}{
\begin{tikzpicture}[very thick] %use thick for all figures!!!
\tikzstyle{every node}=[anchor=west]
%\draw[help lines] (0,0) grid (15,4);
\draw [black] (15,0) -- (15,4) -- (15.5,4) -- (15.5,3.5) -- (21,3.5) -- (21,4) -- (21.5,4) -- (21.5,0) -- (15,0);
\draw [black,fill=black] (16,		0.75) circle [radius=0.05];
\draw [black,fill=black] (16.25,	0.75) circle [radius=0.05];
\draw [black,fill=black] (16.5,		0.75) circle [radius=0.05];
\draw [black,fill=black] (18,		1.5) circle [radius=0.05];
\draw [black,fill=black] (18,		1.75) circle [radius=0.05];
\draw [black,fill=black] (20.5,	2.5) circle [radius=0.05];
\draw [black,fill=black] (17,		2.3) circle [radius=0.05];
\draw [black,fill=black] (17.5,		0.75) circle [radius=0.05];
\draw [black,fill=black] (20,		1) circle [radius=0.05];
\draw [black,fill=black] (19.5,		1.75) circle [radius=0.05];
\draw [black,fill=black] (19,		2.8) circle [radius=0.05];
\draw [black,fill=black] (18.4,		1) circle [radius=0.05];
\draw [black,fill=black] (16.125,	2.4) circle [radius=0.05];

\draw [-latex new, arrow head=3mm,black] (16,3.85) -- (17.75,3.85);
\draw [-latex new, arrow head=3mm,black] (20.5,3.85) -- (18.75,3.85);
\node[anchor=south,black,rotate=0] at (18.25,3.45) {\Large \sffamily pI};

%\node[anchor=south,black,rotate=0] at (1,0.6) {\Large \sffamily pI};
%\draw [arrows=->,black] (4,3) -- (1.5,3);
%\node[anchor=south,black,rotate=0] at (1,0.6) {\Large \sffamily pI};

%\node[anchor=south,black,rotate=0] at (22.25,3) {\Large \sffamily ($-$)};
%\node[anchor=south,black,rotate=0] at (22.25,0) {\Large \sffamily (+)};
%\draw [-latex new, arrow head=3mm,black] (22.25,3) -- (22.25,1);

\draw [-latex new, arrow head=3mm,black] (21,3) -- (21,0.5) node [above,black,midway,rotate=90] {\Large \sffamily size};

\end{tikzpicture}}} \caption{text}
\end{figure}%


\begin{figure}[!htbp]
\centering
\scalebox{0.65}{
\begin{tikzpicture}[very thick] %use thick for all figures!!!
\tikzstyle{every node}=[anchor=west]
%\draw[help lines] (0,0) grid (20,10);

\draw [black] (2,-0.9) -- (2,6.9) -- (1.7,6.9) -- (1.7,-0.9) -- (2,-0.9);
\draw [black] (1.7,-0.9) -- (2,-0.6);
\draw [black] (1.7,-0.6) -- (2,-0.3);
\draw [black] (1.7,-0.3) -- (2,0);
\draw [black] (1.7,0) -- (2,0.3);
\draw [black] (1.7,0.3) -- (2,0.6);
\draw [black] (1.7,0.6) -- (2,0.9);
\draw [black] (1.7,0.9) -- (2,1.2);
\draw [black] (1.7,1.2) -- (2,1.5);
\draw [black] (1.7,1.5) -- (2,1.8);
\draw [black] (1.7,1.8) -- (2,2.1);
\draw [black] (1.7,2.1) -- (2,2.4);
\draw [black] (1.7,2.4) -- (2,2.7);
\draw [black] (1.7,2.7) -- (2,3);
\draw [black] (1.7,3) -- (2,3.3);
\draw [black] (1.7,3.3) -- (2,3.6);
\draw [black] (1.7,3.6) -- (2,3.9);
\draw [black] (1.7,3.9) -- (2,4.2);
\draw [black] (1.7,4.2) -- (2,4.5);
\draw [black] (1.7,4.5) -- (2,4.8);
\draw [black] (1.7,4.8) -- (2,5.1);
\draw [black] (1.7,5.1) -- (2,5.4);
\draw [black] (1.7,5.4) -- (2,5.7);
\draw [black] (1.7,5.7) -- (2,6);
\draw [black] (1.7,6) -- (2,6.3);
\draw [black] (1.7,6.3) -- (2,6.6);
\draw [black] (1.7,6.6) -- (2,6.9);

\draw [black] (2,0.5) to [out=30,in=-90] (3,2.5);
\draw [black] (3,2.5) -- (2.75,2.25) -- (2.5,2.75) -- (2.1,2.5) -- (2.35,3) -- (2.1,3.5) -- (2.5,3.25) -- (2.75,3.75) -- (3,3.5);
\draw [black] (2,5.5) to [out=-30,in=90] (3,3.5);

\draw [black,fill=black] (2.25,1.5) circle [radius=0.075];
\draw [black,fill=black] (2.5,2) circle [radius=0.075];
\draw [black,fill=black] (2.3,2.25) circle [radius=0.075];
\draw [black,fill=black] (2.2,4) circle [radius=0.075];
\draw [black,fill=black] (2.75,4.25) circle [radius=0.075];
\draw [black,fill=black] (2.35,4.75) circle [radius=0.075];

\draw [black,fill=black] (5.2,3.75) circle [radius=0.075];
\draw [black,fill=black] (6.4,2.3) circle [radius=0.075];
\draw [black,fill=black] (4.8,3) circle [radius=0.075];
\draw [black,fill=black] (4,2.5) circle [radius=0.075];
\draw [black,fill=black] (6.75,3.5) circle [radius=0.075];
\draw [black,fill=black] (6,3) circle [radius=0.075];
\draw [black,fill=black] (5.4,2.5) circle [radius=0.075];

\draw [black,ultra thick,rotate=30] (3.75,1.35) -- (10,1.35) -- (10,1.60) -- (12,1.60) -- (12,1.1) -- (10,1.1) -- (10,1.35)  ;
\node[black,rotate=0] at (5,7) {\Large \sffamily laser beam};
\draw [black,dotted] (6,6.5) -- (7,5.6);
\draw [black,-latex new, arrow head=3mm] (3,3) -- (4,3); 
\node[black,rotate=0] at (3,-0.5) {\Large \sffamily matrix/analyte mixture};
\draw [black,dotted] (4,0) -- (2.6,1.75);
\node[black,rotate=0] at (-1,7) {\Large \sffamily target};
\node[black,rotate=0] at (-1,6.26) {\Large \sffamily plate};
\draw [black,dotted] (1,6.5) -- (1.7,5.75);
\draw [black,-latex new, arrow head=3mm] (10.5,3) -- (11.5,3); 
\node[black,rotate=0] at (11.75,3) {\Large \sffamily to mass spectrometer};
\node[black,rotate=0] at (7,1.75) {\Large \sffamily charged};
\node[black,rotate=0] at (7,1) {\Large \sffamily analytes};
\draw [black,dotted] (8,2.25) -- (6.5,3);

\node[black,rotate=0] at (12,7) {\Large \sffamily focussing lenses};
\draw [black,dotted] (13,6.5) -- (11,5);
\draw [black] (10,0) -- (10,2.5) (10,3.5) -- (10,6);
\draw [black] (10.5,0) -- (10.5,2.5) (10.5,3.5) -- (10.5,6);
\draw [black] (11,0) -- (11,2.5) (11,3.5) -- (11,6);

\end{tikzpicture}}
\caption{text}
\label{fig:MALDI}
\end{figure}

\begin{figure}[!htbp]
\centering
\scalebox{0.65}{
\begin{tikzpicture}[very thick] %use thick for all figures!!!
\tikzstyle{every node}=[anchor=west]
%\draw[help lines] (0,0) grid (20,20);

\node[black,rotate=0] at (9,12) {\Large \sffamily flow of liquid analytes (from the RPLC)};
\draw [black,dotted] (8.75,12) -- (5.75,12);

\draw[rotate around={-60:(12,3)},black] (0,2.4) -- (4,2.4) -- (4,1.6) --  (0,1.6)  (4,2.4) -- (5,2) to [out=0,in=205] (10,3.3)  (4,1.6) -- (5,2) to [out=0,in=155] (10,0.7) ;
\draw [rotate around={-60:(12,3)},black,-latex new, arrow head=3mm] (0.5,2) -- (1.5,2); 
\draw [rotate around={-60:(12,3)},black,dashed,-latex new, arrow head=3mm] (3,2.75) -- (4,2.75); 
\draw [rotate around={-60:(12,3)},black,dashed,-latex new, arrow head=3mm] (1.5,1.25) -- (2.5,1.25); 
\draw[rotate around={-60:(12,3)},black,fill=black] (8,2) circle [radius=0.2] (9.75,2) circle [radius=0.1] (9.45,2.2) circle [radius=0.1] (9.5,1.8) circle [radius=0.1]  (11,2) circle [radius=0.05] (11.2,2.2) circle [radius=0.05] (11.4,1.8) circle [radius=0.05] (11.6,2) circle [radius=0.05] (11.8,2.5) circle [radius=0.05] (11.7,1.5) circle [radius=0.05] (11.9,1.7) circle [radius=0.05] (11.5,2.3) circle [radius=0.05] (12,2) circle [radius=0.05];

\draw [black] (13,2.8) -- (14,2.8) -- (14,4) (14.5,2.5) -- (18,2.5);
\draw [black] (13,1.2) -- (14,1.2) -- (14,0) (14.5,1.5) -- (18,1.5);

\node[black,rotate=0] at (15,3.5) {\Large \sffamily drying gas (N$_{\mathsf{2}}$)};
\draw [black,dashed,-latex new, arrow head=3mm] (15,3) to [out=-140,in=0] (13.5,2.5); 
\draw [black,dashed,-latex new, arrow head=3mm] (15,1) to [out=140,in=0] (13.5,1.5); 

\draw [black,-latex new, arrow head=3mm] (12.5,2) -- (14.5,2); 
\node[black,rotate=0] at (15,2) {\Large \sffamily to mass spectrometer };
\node[black,rotate=0] at (6,0) {\Large \sffamily $+$ $-$};
\draw [black] (6.375,0.3) -- (6.375,9.95); 
\draw [black] (7,0.3) -- (7,0.5) -- (13.5,0.5) -- (13.5,1.2); 

\node[black,rotate=0] at (9,11) {\Large \sffamily conducting needle};
\draw [black,dotted] (8.75,11) -- (6.7,11);

\node[black,rotate=0] at (9,10) {\Large \sffamily nebuliser gas (N$_{\mathsf{2}}$)};
\draw [black,dotted] (8.75,10) -- (7.6,10);

\node[black,rotate=0] at (9,9) {\Large \sffamily Taylor cone};
\draw [black,dotted] (8.75,9) -- (7.6,9);
\node[black,rotate=0] at (13,5) {\Large \sffamily spray of droplets};
\draw [black,dotted] (12.75,4.90) -- (10.5,4);
\end{tikzpicture}}
\caption{text}
\label{fig:ESI}
\end{figure}


\begin{figure}[!htbp]
\centering
\scalebox{0.65}{
\begin{tikzpicture}[very thick] %use thick for all figures!!!
\tikzstyle{every node}=[anchor=west]
%\draw[help lines] (-2,0) grid (20,4);

\draw [black] (-1.9,3.9) -- (-0.2,2.1); 
\draw [black] (-2,3.8) -- (-0.3,2); 
\draw [black,-latex new, arrow head=3mm] (0.25,2) -- (1.25,2); 
\node[black,rotate=0] at (-2,0.5) {\Large \sffamily ESI needle};
\draw [black,dotted] (-1.5,1) -- (-0.6,2.25); 

\draw [black] (9,0) -- (9.5,0) to [out=45,in=270] (10.25,1.9) -- (9,1.9) -- (9,0);
\draw [black] (9,4) -- (9.5,4) to [out=-45,in=90] (10.25,2.1) -- (9,2.1) -- (9,4);

\draw [black] (10.25,0) to [out=75,in=180] (11.5,1.5) to [out=0,in=105] (12.75,0) -- (10.25,0);
\draw [black] (10.25,4) to [out=-75,in=180] (11.5,2.5) to [out=0,in=-105] (12.75,4) -- (10.25,4);

\draw [black] (14,0) -- (13.5,0) to [out=135,in=270] (12.75,1.9) -- (14,1.9) -- (14,0);
\draw [black] (14,4) -- (13.5,4) to [out=225,in=90] (12.75,2.1) -- (14,2.1) -- (14,4);

\draw [black] (2,0) -- (2,1.5);
\draw [black] (2,4) -- (2,2.5);
\draw [black] (1.5,1.5) -- (2,1.5);
\draw [black] (1.5,2.5) -- (2,2.5);

\node[black,rotate=0] at (2.7,2.5) {\Large \sffamily ion guide};
%\draw [black,dotted] (4.8,2.75) -- (5.3,3); 
%\draw [black,dotted] (5.1,2.75) -- (6,1);

\draw [black] (3,0) -- (3,0.1);
\draw [black] (3.5,0) -- (3.5,0.25);
\draw [black] (4,0) -- (4,0.5);
\draw [black] (4.5,0) -- (4.5,0.75);
\draw [black] (5,0) -- (5,1);
\draw [black] (5.5,0) -- (5.5,1.25);
\draw [black] (6,0) -- (6,1.5);
\draw [black] (6.5,0) -- (6.5,1.75);
\draw [black] (7,0) -- (7,1.9);
\draw [black] (7.5,0) -- (7.5,1.9);
\draw [black] (8,0) -- (8,1.9);

\draw [black] (3,4) -- (3,3.9);
\draw [black] (3.5,4) -- (3.5,3.75);
\draw [black] (4,4) -- (4,3.5);
\draw [black] (4.5,4) -- (4.5,3.25);
\draw [black] (5,4) -- (5,3);
\draw [black] (5.5,4) -- (5.5,2.75);
\draw [black] (6,4) -- (6,2.5);
\draw [black] (6.5,4) -- (6.5,2.25);
\draw [black] (7,4) -- (7,2.1);
\draw [black] (7.5,4) -- (7.5,2.1);
\draw [black] (8,4) -- (8,2.1);

\node[black,rotate=0] at (15.5,2.3) {\Large \sffamily ions are sequentially};
\node[black,rotate=0] at (15.5,1.7) {\Large \sffamily ejected from the trap};

%\draw [black,arrows=->] (8,2) -- (9,2); 
\draw [black,-latex new, arrow head=3mm] (14,2) -- (15.2,2); 

\node[black,rotate=0] at (15,0.5) {\Large \sffamily end caps};
\draw [black,dotted] (14.7,0.45) -- (9.5,1.5); 
\draw [black,dotted] (14.7,0.45) -- (13.5,1.5);

\node[black,rotate=0] at (15,3.5) {\Large \sffamily ring electrode};
\draw [black,dotted] (14.7,3.5) -- (11.5,3.5); 

\draw [black,fill=black] (12,2) circle [radius=0.08];
\draw [black,fill=black] (11,2) circle [radius=0.08];
\draw [black,fill=black] (11.5,2.2) circle [radius=0.08];
\draw [black,fill=black] (11.7,1.8) circle [radius=0.08];
\draw [black,fill=black] (11.3,1.9) circle [radius=0.08];

\end{tikzpicture}}
\caption{text}
\label{fig:iontrap}
\end{figure}


\begin{figure}[!htbp]
\centering
\renewcommand*\printatom[1]{\ensuremath{\mathsf{#1}}}
\scalebox{1}{
\begin{tikzpicture} 
%\draw[help lines] (0,0) grid (5,2);
\node [] at (0,0) {\setatomsep{0.8cm}\chemfig[black]{H_{2}N-[:0]C(-[:-90]R_{1})(-[:90]H)-[:0]C(=[:90]O)-[:0]NH-[:0]C(-[:-90]R_{2})(-[:90]H)-[:0]C(=[:90]O)-[:0]NH-[:0]C(-[:-90]R_{3})(-[:90]H)(-[:90]H)-[:0]C(=[:90]O)-[:0]NH-[:0]C(-[:-90]R_{4})(-[:90]H)-[:0]C(=[:30]O)-[:-30]OH}};

\draw [arrows=-] [black,thick,dotted] (-3.75,1.4) -- (-3.9,1.4) -- (-3.9,-1.4) -- (-4.05,-1.4);
\draw [arrows=-] [black,thick,dotted] (-2.95,1.4) -- (-3.1,1.4) -- (-3.1,-1.4) -- (-3.25,-1.4) ;
\draw [arrows=-] [black,thick,dotted] (-1.85,1.4) -- (-2,1.4) -- (-2,-1.4) -- (-2.15,-1.4)  ;
\draw [arrows=-] [black,thick,dotted] (-1.05,1.4) -- (-1.2,1.4) -- (-1.2,-1.4) -- (-1.35,-1.4)  ;
\draw [arrows=-] [black,thick,dotted] (-0.25,1.4) -- (-0.4,1.4) -- (-0.4,-1.4) -- (-0.55,-1.4)  ;
\draw [arrows=-] [black,thick,dotted] (0.825,1.4) -- (0.675,1.4) -- (.675,-1.4) -- (0.525,-1.4)  ;
\draw [arrows=-] [black,thick,dotted] (1.625,1.4) -- (1.475,1.4) -- (1.475,-1.4) -- (1.325,-1.4);
\draw [arrows=-] [black,thick,dotted] (2.425,1.4) -- (2.275,1.4) -- (2.275,-1.4) -- (2.125,-1.4);
\draw [arrows=-] [black,thick,dotted] (3.525,1.4) -- (3.375,1.4) -- (3.375,-1.4) -- (3.225,-1.4);

\node[black,anchor=north] at (-3.75,1.95) { \sffamily x$_{\mathsf{3}}$};
\node[black,anchor=north] at (-2.95,1.95) { \sffamily y$_{\mathsf{3}}$};
\node[black,anchor=north] at (-1.85,1.95) { \sffamily z$_{\mathsf{3}}$};

\node[black,anchor=north] at (-1.05,1.95) { \sffamily x$_{\mathsf{2}}$};
\node[black,anchor=north] at (-0.25,1.95) { \sffamily y$_{\mathsf{2}}$};
\node[black,anchor=north] at (0.825,1.95) { \sffamily z$_{\mathsf{2}}$};

\node[black,anchor=north] at (1.625,1.95) { \sffamily x$_{\mathsf{1}}$};
\node[black,anchor=north] at (2.425,1.95) { \sffamily y$_{\mathsf{1}}$};
\node[black,anchor=north] at (3.525,1.95) { \sffamily z$_{\mathsf{1}}$};

\node[black,anchor=south] at (-4.05,-2) { \sffamily a$_{\mathsf{1}}$};
\node[black,anchor=south] at (-3.25,-2) { \sffamily b$_{\mathsf{1}}$};
\node[black,anchor=south] at (-2.15,-2) { \sffamily c$_{\mathsf{1}}$};

\node[black,anchor=south] at (-1.35,-2) { \sffamily a$_{\mathsf{2}}$};
\node[black,anchor=south] at (-0.55,-2) { \sffamily b$_{\mathsf{2}}$};
\node[black,anchor=south] at (0.525,-2) { \sffamily c$_{\mathsf{2}}$};

\node[black,anchor=south] at (1.325,-2) { \sffamily a$_{\mathsf{3}}$};
\node[black,anchor=south] at (2.155,-2) { \sffamily b$_{\mathsf{3}}$};
\node[black,anchor=south] at (3.225,-2) { \sffamily c$_{\mathsf{3}}$};

\draw [black,thick] (-5.75,1.4) -- (-6,1.4) -- (-6,-1.4) -- (-5.75,-1.4);
\draw [black,thick] (5.75,1.4) -- (6,1.4) -- (6,-1.4) -- (5.75,-1.4);
\node[black,anchor=north] at (6.5,1.75) { \sffamily 2H$^{+}$};

\end{tikzpicture} }
\caption{text}
\label{fig:frag}
\end{figure}


\begin{figure}[!htbp]
\centering
\scalebox{0.65}{
\begin{tikzpicture}[] %use thick for all figures!!!
\tikzstyle{every node}=[anchor=west]
%\draw[help lines] (0,0) grid (15,6);

\draw [black,-latex new, arrow head=3mm] (0.8,0) -- (8,0);
\node[black] at (7,-0.35) { \sffamily m/z};
\draw [black,-latex new, arrow head=3mm] (1,-0.2) -- (1,4);
\node[black,rotate=90] at (0.5,2.5) { \sffamily Intensity};
%b-series
%b1
%-
%b2
\draw [black] (1.870713,0) -- (1.870713,2.25);
\node[black,anchor=south] at (1.870713,2.25) { \sffamily \Large b$_{\mathsf{2}}$};
%b3
\draw [black] (3.431724,0) -- (3.431724,1.75);
\node[black,anchor=south] at (3.431724,1.75) { \sffamily \Large b$_{\mathsf{3}}$};
%b4
\draw [black] (5.292518,0) -- (5.292518,3);
\node[black,anchor=south] at (5.292518,3) { \sffamily \Large b$_{\mathsf{4}}$};
%b5
\draw [black] (6.423358,0) -- (6.423358,1.5);
\node[black,anchor=south] at (6.423358,1.5) { \sffamily \Large b$_{\mathsf{5}}$};
%b6
%-
%y-series
%y6
%-
%y5
\draw [black] (6.593624,0) -- (6.593624,1);
\node[black,anchor=south] at (7,1) { \sffamily \Large y$_{\mathsf{5}}$};
%y4
\draw [black] (5.883253,0) -- (5.883253,2.25);
\node[black,anchor=south] at (5.883253,2.25) { \sffamily \Large y$_{\mathsf{4}}$};
%y3
\draw [black] (4.322241,0) -- (4.322241,2.5);
\node[black,anchor=south] at (4.322241,2.5) { \sffamily \Large y$_{\mathsf{3}}$};
%y2
\draw [black] (2.461448,0) -- (2.461448,2.75);
\node[black,anchor=south] at (2.461448,2.75) { \sffamily \Large y$_{\mathsf{2}}$};
%y1
\draw [black] (1.330608,0) -- (1.330608,1.5);
\node[black,anchor=south] at (1.330608,1.5) { \sffamily \Large y$_{\mathsf{1}}$};

\node[black,anchor=west] at (11,4.5) { \sffamily \Large b$_{\mathsf{1}}$\phantom{A}D\phantom{AA}ARWIN\phantom{A}y$_{\mathsf{5}}$};
\node[black,anchor=west] at (11,3.5) { \sffamily \Large b$_{\mathsf{2}}$\phantom{A}DA\phantom{AA}RWIN\phantom{A}y$_{\mathsf{4}}$};
\node[black,anchor=west] at (11,2.5) { \sffamily \Large b$_{\mathsf{3}}$\phantom{A}DAR\phantom{AA}WIN\phantom{A}y$_{\mathsf{3}}$};
\node[black,anchor=west] at (11,1.5) { \sffamily \Large b$_{\mathsf{4}}$\phantom{A}DARW\phantom{AA}IN\phantom{A}y$_{\mathsf{2}}$};
\node[black,anchor=west] at (11,0.5) { \sffamily \Large b$_{\mathsf{5}}$\phantom{A}DARWI\phantom{AA}N\phantom{A}y$_{\mathsf{1}}$};

\draw [black,latex new-latex new, arrow head=3mm] (2.461448,4) -- (4.322241,4);
\draw [black,latex new-latex new, arrow head=3mm] (3.431724,4.5) -- (5.292518,4.5);

\end{tikzpicture}}
\caption{text}
\label{fig:denovo}
\end{figure}


\begin{figure}[!htbp]
\centering
\scalebox{0.4}{
\begin{tikzpicture}[very thick] %use thick for all figures!!!
\tikzstyle{every node}=[anchor=west]
%\draw[help lines] (0,0) grid (20,15);

\draw [black,-latex new, arrow head=3mm] (0.8,5) -- (8,5);
%\node[black] at (7,5.35) { \sffamily m/z};
\draw [black,-latex new, arrow head=3mm] (1,4.8) -- (1,8);
%\node[black,rotate=90] at (0.5,6.5) { \sffamily Intensity};

\draw [black] (1.870713,5) -- (1.870713,7.25);
\draw [black] (3.431724,5) -- (3.431724,6.75);
\draw [black] (5.292518,5) -- (5.292518,8);
\draw [black] (6.423358,5) -- (6.423358,6.5);
\draw [black] (6.593624,5) -- (6.593624,6);
\draw [black] (5.883253,5) -- (5.883253,7.25);
\draw [black] (4.322241,5) -- (4.322241,7.5);
\draw [black] (2.461448,5) -- (2.461448,7.75);
\draw [black] (1.330608,5) -- (1.330608,6.5);

%DARWIN
\draw [black,-latex new, arrow head=3mm] (15.8,10) -- (23,10);
%\node[black] at (22,10.35) { \sffamily m/z};
\draw [black,-latex new, arrow head=3mm] (16,9.8) -- (16,13);
%\node[black,rotate=90] at (15.5,11.4) { \sffamily Intensity};

\draw [black] (16.870713,10) -- (16.870713,12.25);
\draw [black] (18.431724,10) -- (18.431724,11.75);
\draw [black] (20.292518,10) -- (20.292518,13);
\draw [black] (21.423358,10) -- (21.423358,11.5);
\draw [black] (21.593624,10) -- (21.593624,11);
\draw [black] (20.883253,10) -- (20.883253,12.25);
\draw [black] (19.322241,10) -- (19.322241,12.5);
\draw [black] (17.461448,10) -- (17.461448,12.75);
\draw [black] (16.330608,10) -- (16.330608,11.5);
\node[black,anchor=west] at (17.5,13.75) { \sffamily \Huge DARWIN};
\draw [black,dotted,-latex new, arrow head=6mm] (23.5,11) -- (26.5,11);
\node[black,anchor=west] at (27,11.25) { \sffamily \Huge best match};

%ARCHIMEDES
\draw [black,-latex new, arrow head=3mm] (15.8,5) -- (23,5);
%\node[black] at (22,5.35) { \sffamily m/z};
\draw [black,-latex new, arrow head=3mm] (16,4.8) -- (16,8);
%\node[black,rotate=90] at (15.5,6.5) { \sffamily Intensity};

\draw [black] (16.5,5) -- (16.5,7.25);
\draw [black] (18,5) -- (18,6.75);
\draw [black] (20,5) -- (20,7.5);
\draw [black] (21,5) -- (21,6.5);
\draw [black] (21.5,5) -- (21.5,6);
\draw [black] (19.5,5) -- (19.5,7.25);
\draw [black] (19,5) -- (19,7.5);
\draw [black] (17,5) -- (17,7.75);
\draw [black] (16.25,5) -- (16.25,6.5);
\node[black,anchor=west] at (16.5,8.55) { \sffamily \Huge ARCHIMEDES};

%EINSTEIN
\draw [black,-latex new, arrow head=3mm] (15.8,0) -- (23,0);
%\node[black] at (22,0.35) { \sffamily m/z};
\draw [black,-latex new, arrow head=3mm] (16,-0.2) -- (16,3);
%\node[black,rotate=90] at (15.5,1.5) { \sffamily Intensity};

\draw [black] (16.75,0) -- (16.75,2.25);
\draw [black] (18,0) -- (18,1);
\draw [black] (20,0) -- (20,2);
\draw [black] (21,0) -- (21,2);
\draw [black] (21.25,0) -- (21.25,1);
\draw [black] (20.75,0) -- (20.75,1.25);
\draw [black] (18.5,0) -- (18.5,2.5);
\draw [black] (17,0) -- (17,2.75);
\draw [black] (16.5,0) -- (16.5,1.5);
\node[black,anchor=west] at (17.5,3.5) { \sffamily \Huge EINSTEIN};

\draw [black,latex new-latex new, arrow head=6mm] (8.5,6.5) -- (13.5,11);
\draw [black,latex new-latex new, arrow head=6mm] (8.5,6) -- (13.5,6);
\draw [black,latex new-latex new, arrow head=6mm] (8.5,5.5) -- (13.5,1);

\draw [black] (19.25,15) ellipse (5cm and 0.5cm);
\draw [black] (14.25,0) -- (14.25,15);
\draw [black] (24.25,0) -- (24.25,15);
\draw[black] (17.25,0) +(180:3) arc (180:360:5 and 0.5) ;

\end{tikzpicture}}
\caption{text}\label{fig:spectral}
\end{figure}

\begin{figure}[!htbp]
\centering
\scalebox{0.4}{
\begin{tikzpicture}[very thick] %use thick for all figures!!!
\tikzstyle{every node}=[anchor=west]
%\draw[help lines] (0,0) grid (20,15);

\draw [black,-latex new, arrow head=3mm] (0.8,5) -- (8,5);
%\node[black] at (7,5.35) { \sffamily m/z};
\draw [black,-latex new, arrow head=3mm] (1,4.8) -- (1,8);
%\node[black,rotate=90] at (0.5,6.5) { \sffamily Intensity};

\draw [black] (1.870713,5) -- (1.870713,7.25);
\draw [black] (3.431724,5) -- (3.431724,6.75);
\draw [black] (5.292518,5) -- (5.292518,8);
\draw [black] (6.423358,5) -- (6.423358,6.5);
\draw [black] (6.593624,5) -- (6.593624,6);
\draw [black] (5.883253,5) -- (5.883253,7.25);
\draw [black] (4.322241,5) -- (4.322241,7.5);
\draw [black] (2.461448,5) -- (2.461448,7.75);
\draw [black] (1.330608,5) -- (1.330608,6.5);

%DARWIN
\draw [black,-latex new, arrow head=3mm] (15.8,10) -- (23,10);
%\node[black] at (22,10.35) { \sffamily m/z};
\draw [black,-latex new, arrow head=3mm] (16,9.8) -- (16,13);
%\node[black,rotate=90] at (15.5,11.4) { \sffamily Intensity};

\draw [black] (16.870713,10) -- (16.870713,12.5);
\draw [black] (18.431724,10) -- (18.431724,12.5);
\draw [black] (20.292518,10) -- (20.292518,12.5);
\draw [black] (21.423358,10) -- (21.423358,12.5);
\draw [black] (21.593624,10) -- (21.593624,12.5);
\draw [black] (20.883253,10) -- (20.883253,12.5);
\draw [black] (19.322241,10) -- (19.322241,12.5);
\draw [black] (17.461448,10) -- (17.461448,12.5);
\draw [black] (16.330608,10) -- (16.330608,12.5);

%ARCHIMEDES
\draw [black,-latex new, arrow head=3mm] (15.8,5) -- (23,5);
%\node[black] at (22,5.35) { \sffamily m/z};
\draw [black,-latex new, arrow head=3mm] (16,4.8) -- (16,8);
%\node[black,rotate=90] at (15.5,6.5) { \sffamily Intensity};

\draw [black] (16.5,5) -- (16.5,7.5);
\draw [black] (18,5) -- (18,7.5);
\draw [black] (20,5) -- (20,7.5);
\draw [black] (21,5) -- (21,7.5);
\draw [black] (21.5,5) -- (21.5,7.5);
\draw [black] (19.5,5) -- (19.5,7.5);
\draw [black] (19,5) -- (19,7.5);
\draw [black] (17,5) -- (17,7.5);
\draw [black] (16.25,5) -- (16.25,7.5);

%EINSTEIN
\draw [black,-latex new, arrow head=3mm] (15.8,0) -- (23,0);
%\node[black] at (22,0.35) { \sffamily m/z};
\draw [black,-latex new, arrow head=3mm] (16,-0.2) -- (16,3);
%\node[black,rotate=90] at (15.5,1.5) { \sffamily Intensity};

\draw [black] (16.75,0) -- (16.75,2.5);
\draw [black] (18,0) -- (18,2.5);
\draw [black] (20,0) -- (20,2.5);
\draw [black] (21,0) -- (21,2.5);
\draw [black] (21.25,0) -- (21.25,2.5);
\draw [black] (20.75,0) -- (20.75,2.5);
\draw [black] (18.5,0) -- (18.5,2.5);
\draw [black] (17,0) -- (17,2.5);
\draw [black] (16.5,0) -- (16.5,2.5);

%\draw [black,dotted,arrows=<->] (8.5,6.5) -- (13.5,11);
\draw [black,latex new-latex new, arrow head=6mm] (8.5,6) -- (13.5,6);
\draw [black,latex new-latex new, arrow head=6mm] (8.5,5.5) -- (13.5,1);

\draw [black,latex new-latex new, arrow head=6mm] (8.5,6.5) -- (13.5,11) node [left,near end,text width=4cm,text height=0.8cm] {\sffamily \Huge best match
\\
\phantom{text}};

\draw [black] (32.25,8.5) ellipse (4cm and 0.5cm);
\draw [black] (28.25,3.5) -- (28.25,8.5);
\draw [black] (36.25,3.5) -- (36.25,8.5);
\draw[black] (31.25,3.5) +(180:3) arc (180:360:4 and 0.5) ;

\node[black,anchor=west] at (30,7) { \sffamily \Huge DARWIN};
\node[black,anchor=west] at (30,6) { \sffamily \Huge ARCHIMEDES};
\node[black,anchor=west] at (30,5) { \sffamily \Huge EINSTEIN};
\node[black,anchor=west] at (30,4) { \sffamily \Huge ...};

\draw [black,-latex new, arrow head=6mm] (28,7) -- (24,11);
\draw [black,-latex new, arrow head=6mm] (28,6) -- (24,6);
\draw [black,-latex new, arrow head=6mm] (28,5) -- (24,1);
\end{tikzpicture}}
\caption{text}
\label{fig:proteindatabase}
\end{figure}


\begin{figure}[!htbp]
\centering
\scalebox{0.4}{
\begin{tikzpicture}[very thick] %use thick for all figures!!!
\tikzstyle{every node}=[anchor=west]
%\draw[help lines] (0,0) grid (20,15);

\draw [black,arrows=->] (0.8,5) -- (8,5);
%\node[black] at (7,5.35) { \sffamily m/z};
\draw [black,arrows=->] (1,4.8) -- (1,8);
%\node[black,rotate=90] at (0.5,6.5) { \sffamily Intensity};

\draw [black] (1.870713,5) -- (1.870713,7.25);
\draw [black] (3.431724,5) -- (3.431724,6.75);
\draw [black] (5.292518,5) -- (5.292518,8);
\draw [black] (6.423358,5) -- (6.423358,6.5);
\draw [black] (6.593624,5) -- (6.593624,6);
\draw [black] (5.883253,5) -- (5.883253,7.25);
\draw [black] (4.322241,5) -- (4.322241,7.5);
\draw [black] (2.461448,5) -- (2.461448,7.75);
\draw [black] (1.330608,5) -- (1.330608,6.5);


\draw [black] (17.25,10.5) ellipse (4cm and 0.5cm);
\draw [black] (13.25,2) -- (13.25,10.5);
\draw [black] (21.25,2) -- (21.25,10.5);
\draw[black] (16.25,2) +(180:3) arc (180:360:4 and 0.5) ;
\draw[black] (16.25,6.1) +(180:3) arc (180:360:4 and 0.5) ;

\node[black,anchor=west] at (15,9) { \sffamily \Huge DARWIN};
\node[black,anchor=west] at (15,8) { \sffamily \Huge ARCHIMEDES};
\node[black,anchor=west] at (15,7) { \sffamily \Huge EINSTEIN};
\node[black,anchor=west] at (15,6) { \sffamily \Huge ...};

\node[black,anchor=west] at (15,5) { \sffamily \Huge NIWRAD};
\node[black,anchor=west] at (15,4) { \sffamily \Huge SEDEMIHCRA};
\node[black,anchor=west] at (15,3) { \sffamily \Huge NIETSNIE};
\node[black,anchor=west] at (15,2) { \sffamily \Huge ...};

\node[black,anchor=west] at (25,10) { \sffamily \Huge PSM};
\node[black,anchor=west] at (25,9) { \sffamily \Huge DARWIN};
\node[black,anchor=west] at (25,8) { \sffamily \Huge ARCHIMEDES};
\node[black,anchor=west] at (25,7) { \sffamily \Huge NIWRAD};
\node[black,anchor=west] at (25,6) { \sffamily \Huge EINSTEIN};
\node[black,anchor=west] at (25,5) { \sffamily \Huge SEDEMIHCRA};
\node[black,anchor=west] at (25,4) { \sffamily \Huge NIETSNIE};
\node[black,anchor=west] at (25,3) { \sffamily \Huge ...};

\node[black,anchor=center] at (32.5,10) { \sffamily \Huge Score};
\node[black,anchor=east] at (33,9) { \sffamily \Huge 10};
\node[black,anchor=east] at (33,8) { \sffamily \Huge 8};
\node[black,anchor=east] at (33,7) { \sffamily \Huge 6};
\node[black,anchor=east] at (33,6) { \sffamily \Huge 4};
\node[black,anchor=east] at (33,5) { \sffamily \Huge 2};
\node[black,anchor=east] at (33,4) { \sffamily \Huge 1};
\node[black,anchor=east] at (33,3) { \sffamily \Huge ...};

\draw [black,-latex new, arrow head=6mm] (21.5,6.1) -- (24,6.1);
\draw [black,-latex new, arrow head=6mm] (8.5,6.1) -- (12,4);
\draw [black,-latex new, arrow head=6mm] (8.5,6.1) -- (12,8) ;

\draw [black] (24.5,9.5) -- (34,9.5);
\draw [black,dashed] (24.5,5.5) -- (34,5.5);
\node[black,anchor=west] at (34.75,5.5) { \sffamily \Huge $T_{s}$};

%\node[black,anchor=west] at (25,1.5) { \sffamily \Huge FDR = ${N_{decoy\ T_{s}}}/{N_{target\ T_{s}}}$};

\end{tikzpicture}}
\scalebox{0.65}{
\begin{tikzpicture}
%\draw [black,-latex new, arrow head=6mm] (7,3) -- (8,4);
\node[black,anchor=west] at (-13,1) { \sffamily \Large };
\node[black,anchor=north west] at (1,1.5) { \sffamily \Large FDR = ${N_{decoy}\ (T_{s})}/{N_{target}\ (T_{s})}$};

\end{tikzpicture}}
\caption{text}
\label{fig:targetdecoy}
\end{figure}%


\begin{figure}[!htbp]
\centering
\scalebox{0.4}{
\begin{tikzpicture}[very thick] %use thick for all figures!!!
\tikzstyle{every node}=[anchor=west]
%\draw[help lines] (0,0) grid (20,15);

\draw [black,arrows=->] (0.8,5) -- (8,5);
%\node[black] at (7,5.35) { \sffamily m/z};
\draw [black,arrows=->] (1,4.8) -- (1,8);
%\node[black,rotate=90] at (0.5,6.5) { \sffamily Intensity};

\draw [black] (1.870713,5) -- (1.870713,7.25);
\draw [black] (3.431724,5) -- (3.431724,6.75);
\draw [black] (5.292518,5) -- (5.292518,8);
\draw [black] (6.423358,5) -- (6.423358,6.5);
\draw [black] (6.593624,5) -- (6.593624,6);
\draw [black] (5.883253,5) -- (5.883253,7.25);
\draw [black] (4.322241,5) -- (4.322241,7.5);
\draw [black] (2.461448,5) -- (2.461448,7.75);
\draw [black] (1.330608,5) -- (1.330608,6.5);

\draw [black] (17.25,10.5) ellipse (4cm and 0.5cm);
\draw [black,dotted] (13.25,2) -- (13.25,6.1);
\draw [black] (13.25,6.1) -- (13.25,10.5);
\draw [black,dotted] (21.25,2) -- (21.25,6.1);
\draw [black] (21.25,6.1) -- (21.25,10.5);
\draw[black,dotted] (16.25,2) +(180:3) arc (180:360:4 and 0.5) ;
\draw[black] (16.25,6.1) +(180:3) arc (180:360:4 and 0.5) ;

\node[black,anchor=west] at (15,9) { \sffamily \Huge DARWIN};
\node[black,anchor=west] at (15,8) { \sffamily \Huge ARCHIMEDES};
\node[black,anchor=west] at (15,7) { \sffamily \Huge EINSTEIN};
\node[black,anchor=west] at (15,6) { \sffamily \Huge ...};

\node[black,anchor=west] at (15,5) { \sffamily \Huge NIWRAD};
\node[black,anchor=west] at (15,4) { \sffamily \Huge SEDEMIHCRA};
\node[black,anchor=west] at (15,3) { \sffamily \Huge NEITSNIE};
\node[black,anchor=west] at (15,2) { \sffamily \Huge ...};

\draw [dotted,black,-latex new, arrow head=6mm] (8.5,6.1) -- (12,4);
\draw [black,-latex new, arrow head=6mm] (8.5,6.1) -- (12,8) ;
\draw [black,-latex new, arrow head=6mm] (21.5,6.1) -- (24,6.1);

\node[black,anchor=west] at (13,1) { \sffamily \large };

\end{tikzpicture}} 
%\hspace{0.3cm}
\scalebox{0.65}{
\begin{tikzpicture}[
    declare function={gamma(\z)=
    2.506628274631*sqrt(1/\z)+ 0.20888568*(1/\z)^(1.5)+ 0.00870357*(1/\z)^(2.5)- (174.2106599*(1/\z)^(3.5))/25920- (715.6423511*(1/\z)^(4.5))/1244160)*exp((-ln(1/\z)-1)*\z;},
    declare function={gammapdf(\x,\k,\theta) = 1/(\theta^\k)*1/(gamma(\k))*\x^(\k-1)*exp(-\x/\theta);}
]

\newcommand\gauss[2]{1/(#2*sqrt(2*pi))*exp(-((x-#1)^2)/(2*#2^2))} 

%\draw[help lines] (0,0) grid (20,5);
 
%X-axis
\draw [very thick, black,-latex new, arrow head=3mm] (-0.2,0.46) -- (7,0.46);
\node[black,anchor=west] at (6,0) { \sffamily \large S};

%Y-axis
\draw [very thick, black,-latex new, arrow head=3mm] (0,0.26) -- (0,6);
\node[black,anchor=west,rotate=90] at (-0.375,3) { \sffamily \large \# spectra};

\node[black,anchor=west] at (2,5) { \sffamily \Large incorrect PSMs};
\draw [very thick, black,dotted] (1.2,4) -- (2,4.75);
\node[black,anchor=west] at (2,3) { \sffamily \Large correct PSMs};
\draw [very thick, black,dotted] (4.5,1) -- (5,2.8);

\draw [very thick, black] (0.4,0.46) -- (0.54,3) -- (0.65,2.8) -- (0.9,5.5) -- (1,5.6) -- (1.4,4) -- (1.5,3) -- (1.6,2) -- (1.8,1) -- (2.2,1.4) -- (2.4,1.3) -- (2.6,1.5) -- (2.7,1.4) -- (3,1.8) -- (3.3,2) -- (3.6,1.7) -- (3.8,1.8) -- (4,1.6) -- (4.2,1.7) -- (4.5,1.3) -- (5,1) -- (6,0.46);


%\draw [black] (2,0.4) to [out=30,in=150] (6,0.4);

\begin{axis}[axis lines=none,samples=50]
\addplot [thick,dashed,smooth, domain=0:40, black] {gammapdf(x,5,0.75)};
\addplot[thick,dashed,smooth, domain=0:40, black] {\gauss{20}{6}};
\end{axis}
\end{tikzpicture}}

\scalebox{0.65}{
\begin{tikzpicture}
%\draw [black,-latex new, arrow head=6mm] (7,3) -- (8,4);
\node[black,anchor=west] at (-5,4) { \sffamily \Large };
\node[black,anchor=north west] at (1,3.5) { \sffamily \Large FDR = $N_{inc}(P_{T})/N(P_{T})$};
\node[black,anchor=north west] at (8.5,3.5) { \sffamily \Large $N_{inc}(P_{T})$ = $\sum\limits_{{x|P_{x}\geq P_{T}}}(1 - P_{x}$)};
\node[black,anchor=north west] at (8.5,2) { \sffamily \Large $N(P_{T})$ = $\sum\limits_{{x|P_{x}\geq P_{T}}}1$};

\end{tikzpicture}}

\caption{text}
\label{fig:peptideprophet}
\end{figure}

\begin{figure}[!htbp]
\centering
\scalebox{0.9}{
\begin{tikzpicture}[very thick] %use thick for all figures!!!
\tikzstyle{every node}=[anchor=west]
%\draw[help lines] (0,0) grid (15,15);

\node[black] at (0,4) { \sffamily Protein A};
\node[black] at (0,3.5) { \sffamily Protein B};

\node[black] at (2,4.5) { \sffamily Peptide 1};
\node[black] at (4,4.5) { \sffamily Peptide 2};
\node[black] at (6,4.5) { \sffamily Peptide 3};

\draw [black,fill=black] (3,4) circle [radius=0.1];
\draw [black,fill=black] (5,4) circle [radius=0.1];
\draw [black,fill=black] (7,4) circle [radius=0.1];
\draw [black,fill=black] (5,3.5) circle [radius=0.1];
\draw [black,fill=black] (7,3.5) circle [radius=0.1];

\node[black] at (10.5,6) { \sffamily 1 a. Prob Protein A = 1};
\node[black] at (11,5.5) { \sffamily Peptide 1 $\rightarrow$ w = 1};
\node[black] at (11,5) { \sffamily Peptide 2 $\rightarrow$ w = 1 };
\node[black] at (11,4.5) { \sffamily Peptide 3 $\rightarrow$ w = 1 };

\node[black] at (10.5,3.5) { \sffamily 1 b. Prob Protein B = 0};
\node[black] at (11,3) { \sffamily Peptide 2 $\rightarrow$ w = 0 };
\node[black] at (11,2.5) { \sffamily Peptide 3 $\rightarrow$ w = 0 };

\end{tikzpicture}}
\caption{text}
\label{fig:weight}
\end{figure}


%%%%%%%%%%%%%%%%%% Objectives %%%%%%%%%%%%%%%%



\newpage{\thispagestyle{empty}\cleardoublepage}
\chapter{Rationale and Objectives} \label{chap:objectives}
\lettergroup{\thechapter}

text
%%%%%%%%%%%%%%%%%% FDR %%%%%%%%%%%%%%%%



\newpage{\thispagestyle{empty}\cleardoublepage}
\chapter{Partially Sequenced Organisms, Decoy Searches and False Discovery Rates} 
\label{chap:FDR}\lettergroup{\thechapter}

\begin{quote}
\textbf{Abstract:} text

\end{quote}

\vfill

Adapted from:\\
\textbf{Victor, B.}, Gabri\"el, S., Kanobana, K., Mostovenko, E., Polman, K., Dorny, P., Deelder, A. M., and Palmblad, M. Partially Sequenced Organisms, Decoy Searches and False Discovery Rates. \textit{Journal of Proteome Research} 11, 3 (2012), 1991--1995.
\newpage

text

\newpage{\thispagestyle{empty}\cleardoublepage}
\chapter{Proteomic Analysis of \textit{Taenia solium} Metacestode Excretion/Secretion Proteins}
\label{chap:ESP} \lettergroup{\thechapter}

\begin{quote}
\textbf{Abstract:} text
\end{quote}

\vfill

Adapted from:\\
\textbf{Victor, B.}, Kanobana, K., Gabri\"el, S., Polman, K., Deckers, N., Dorny, P., Deelder, A. M., and Palmblad, M. Proteomic Analysis of \textit{Taenia solium} Metacestode Excretion/Secretion Proteins. \textit{Proteomics} 12, 11 (2012), 1860--1869.
\newpage

text

\newpage{\thispagestyle{empty}\cleardoublepage}
\chapter{Use of Expressed Sequence Tags for the Identification of \textit{Taenia solium} Metacestode Excretion/Secretion Proteins} 
\chaptermark{Use of ESTs for the Identification of \textit{T. solium} Metacestode ESPs.}
\label{chap:EST} \lettergroup{\thechapter}

\begin{quote}
\textbf{Abstract:} text
\end{quote}

\vfill

Adapted from:\\
\textbf{Victor, B.}, Dorny, P., Kanobana, K., Polman, K., Lindh, J., Deelder, A. M., Palmblad, M., and Gabri\"el, S. Use of Expressed Sequence Tags as an Alternative Approach for the Identification of \textit{Taenia solium} Metacestode Excretion/Secretion Proteins. \textit{BMC Research Notes} 6, 1 (2013), 224.
\newpage

text

\newpage{\thispagestyle{empty}\cleardoublepage}
\chapter{The Search for Differential Proteins}  
\label{chap:diffproteins}\lettergroup{\thechapter}

\begin{figure}[!htbp]
\centering
{\sffamily
\begin{tikzpicture} 
\tikzstyle{every node}=[anchor=center]
\tikzstyle{column 3}=[text width=0.3\linewidth,text centered]
\tikzstyle{column 2}=[text width=0.3\linewidth,text centered]
\tikzstyle{column 1}=[text width=0.3\linewidth,text centered]
\matrix [row sep=0.2cm,column sep=0.5cm] %draw 
{
\node()[]{}; 	&	\node(vc)[]{tandem mass spectro-\\metry analysis of \textit{Taenia} excretion/secretion proteins}; &\node()[]{}; \\ 

\node(v1)[]{all \textit{Taenia}\\ \textit{solium} spectra}; 	&	&\node(v3)[]{all \textit{Taenia}\\ \textit{hydatigena} spectra}; \\ 

 &  \node(v5)[]{host proteins and \\common contaminations};  		&\\ 

\node(v7){spectral\\ quality filter};& & \node(v9)[]{spectral\\ quality filter}; 	\\ 

	& \node(v11)[]{non-peptide and  \\low quality spectra}; 			&\\ 
\node(v13){high quality \textit{Taenia}\\ \textit{solium} spectra};	&  \node(v14){}; 		& \node(v15){high quality \textit{Taenia}\\ \textit{hydatigena} spectra}; \\

\node(){};	&  \node(va){spectra are binned prior to spectral comparison}; 		& \node(){}; \\

\node(v16){spectrum 1: \centering
\scalebox{0.25}{
\begin{tikzpicture}[very thick] %use thick for all figures!!!
\tikzstyle{every node}=[anchor=west]
\draw [black,-latex new, arrow head=3mm] (0.8,5) -- (8,5);
\draw [black,-latex new, arrow head=3mm] (1,4.8) -- (1,8);
\draw [black] (1.5,5) rectangle (2,6);
\draw [black] (2,5) rectangle (2.5,6.5);
\draw [black] (2.5,5) rectangle (3,6.5);
\draw [black] (3,5) rectangle (3.5,8);
\draw [black] (3.5,5) rectangle (4,7);
\draw [black] (4,5) rectangle (4.5,6.5);
\draw [black] (4.5,5) rectangle (5,7.5);
\draw [black] (5,5) rectangle (5.5,7);
\draw [black] (5.5,5) rectangle (6,5.5);
\draw [black] (6,5) rectangle (6.5,6);
\draw [black] (6.5,5) rectangle (7,5.5);

\end{tikzpicture}}};	&  \node(v17){}; 		& \node(v18){spectrum 1: \centering
\scalebox{0.25}{
\begin{tikzpicture}[very thick] %use thick for all figures!!!
\tikzstyle{every node}=[anchor=west]
\draw [black,-latex new, arrow head=3mm] (0.8,5) -- (8,5);
\draw [black,-latex new, arrow head=3mm] (1,4.8) -- (1,8);
\draw [black] (1.5,5) rectangle (2,6);
\draw [black] (2,5) rectangle (2.5,6.5);
\draw [black] (2.5,5) rectangle (3,6.5);
\draw [black] (3,5) rectangle (3.5,8);
\draw [black] (3.5,5) rectangle (4,7);
\draw [black] (4,5) rectangle (4.5,6.5);
\draw [black] (4.5,5) rectangle (5,7.5);
\draw [black] (5,5) rectangle (5.5,7);
\draw [black] (5.5,5) rectangle (6,5.5);
\draw [black] (6,5) rectangle (6.5,6);
\draw [black] (6.5,5) rectangle (7,5.5);

\end{tikzpicture}}}; \\

\node(v19){spectrum 2: \centering
\scalebox{0.25}{
\begin{tikzpicture}[very thick] %use thick for all figures!!!
\tikzstyle{every node}=[anchor=west]

\draw [black,-latex new, arrow head=3mm] (0.8,5) -- (8,5);
\draw [black,-latex new, arrow head=3mm] (1,4.8) -- (1,8);
\draw [black] (1.5,5) rectangle (2,6);
\draw [black] (2,5) rectangle (2.5,7.5);
\draw [black] (2.5,5) rectangle (3,7);
\draw [black] (3,5) rectangle (3.5,7.25);
\draw [black] (3.5,5) rectangle (4,7);
\draw [black] (4,5) rectangle (4.5,6);
\draw [black] (4.5,5) rectangle (5,7.5);
\draw [black] (5,5) rectangle (5.5,5.5);
\draw [black] (5.5,5) rectangle (6,5.75);
\draw [black] (6,5) rectangle (6.5,6);
\draw [black] (6.5,5) rectangle (7,5.5);

\end{tikzpicture}}};	&  \node(v20){}; 		& \node(v21){spectrum 2: \centering
\scalebox{0.25}{
\begin{tikzpicture}[very thick] %use thick for all figures!!!
\tikzstyle{every node}=[anchor=west]

\draw [black,-latex new, arrow head=3mm] (0.8,5) -- (8,5);
\draw [black,-latex new, arrow head=3mm] (1,4.8) -- (1,8);
\draw [black] (1.5,5) rectangle (2,6);
\draw [black] (2,5) rectangle (2.5,6.5);
\draw [black] (2.5,5) rectangle (3,7.5);
\draw [black] (3,5) rectangle (3.5,8);
\draw [black] (3.5,5) rectangle (4,7);
\draw [black] (4,5) rectangle (4.5,6.5);
\draw [black] (4.5,5) rectangle (5,5.5);
\draw [black] (5,5) rectangle (5.5,6);
\draw [black] (5.5,5) rectangle (6,8);
\draw [black] (6,5) rectangle (6.5,6);
\draw [black] (6.5,5) rectangle (7,5.5);

\end{tikzpicture}}}; \\

\node(v25){spectrum 3: ...\phantom{\centering
\scalebox{0.25}{
\begin{tikzpicture}[very thick] %use thick for all figures!!!
\tikzstyle{every node}=[anchor=west]
\draw [black,-latex new, arrow head=3mm] (0.8,5) -- (6.5,5);
\draw [black,-latex new, arrow head=3mm] (1,4.8) -- (1,8);
\end{tikzpicture}}}};	&  \node(v26){}; 		& \node(v27){spectrum 3: ...\phantom{\centering
\scalebox{0.25}{
\begin{tikzpicture}[very thick] %use thick for all figures!!!
\tikzstyle{every node}=[anchor=west]
\draw [black,-latex new, arrow head=3mm] (0.8,5) -- (6.5,5);
\draw [black,-latex new, arrow head=3mm] (1,4.8) -- (1,8);
\end{tikzpicture}}}}; \\

\node(){};	&  \node(){\phantom{text}\\ \phantom{text}}; 		& \node(){}; \\

\node(v28){unique \textit{Taenia}\\ \textit{solium} spectra};	&  \node(vb){six reading frame translation of the \\\textit{Taenia solium} genome};		& \node(v30){unique \textit{Taenia}\\ \textit{hydatigena} spectra}; \\

\node(){};	&  \node(){all identified peptides +\\ their matching spectra}; 		& \node(){}; \\

\node(v31){unique \textit{Taenia solium} peptides (proteins)};	&  \node(v32){}; 		& \node(v33){}; \\
}; 

\draw [-latex new, arrow head=2mm] [black,thick] (vc.west) to [out=180,in=45] (v1.north);
\draw [-latex new, arrow head=2mm] [black,thick] (vc.east) to [out=0,in=135] (v3.north);

\draw [-latex new, arrow head=2mm] [black,thick] (v1.south) -- (v7.north);
\draw [-latex new, arrow head=2mm] [black,thick] (v7.south) -- (v13.north);
\draw [-latex new, arrow head=2mm] [black,thick] (v13.south) -- (v16.north);
\draw [-latex new, arrow head=2mm] [black,thick] (v25.south) -- (v28.north);
\draw [-latex new, arrow head=2mm] [black,thick] (v28.south) -- (v31.north);

\draw [-latex new, arrow head=2mm] [black,thick] (v3.south) -- (v9.north);
\draw [-latex new, arrow head=2mm] [black,thick] (v9.south) -- (v15.north);
\draw [-latex new, arrow head=2mm] [black,thick] (v15.south) -- (v18.north);
\draw [-latex new, arrow head=2mm] [black,thick] (v27.south) -- (v30.north);

\draw [-latex new, arrow head=2mm] [black,thick] (v1.south) to [out=-45,in=180] (v5.west);
\draw [-latex new, arrow head=2mm] [black,thick] (v7.south) to [out=-45,in=180] (v11.west);

\draw [-latex new, arrow head=2mm] [dotted,black,thick] (v1.south) to [out=-45,in=125] (vb.155);
\draw [-latex new, arrow head=2mm] [dotted,black,thick] (vb.-155) to [out=-125,in=90] (v31.north);

\draw [-latex new, arrow head=2mm] [black,thick] (va.west) to [out=180,in=45] (v16.20);
\draw [-latex new, arrow head=2mm] [black,thick] (va.east) to [out=0,in=135] (v18.160);

%\draw [-latex new, arrow head=2mm] [black,thick] (vb.west) to [out=180,in=45] (v31.north);

\draw [-latex new, arrow head=2mm] [black,thick] (v3.south) to [out=-135,in=0] (v5.east);
\draw [-latex new, arrow head=2mm] [black,thick] (v9.south) to [out=-135,in=0] (v11.east);

\draw [latex-latex new, arrow head=2mm,dashed] [black,thick] (v16.east) -- (v18.west);
\draw [latex-latex new, arrow head=2mm,dashed] [black,thick] (v16.east) -- (v21.west);
%\draw [latex-latex new, arrow head=2mm,dashed] [black,thick] (v16.east) -- (v24.west);
\draw [latex-latex new, arrow head=2mm,dashed] [black,thick] (v16.east) -- (v27.west) node[left,midway,text width=1.8cm,text height=1.2cm]{\textcolor{black}{spectral \\comparisons}};

%\draw [-latex new, arrow head=2mm] [black,thick] (v4.south) -- (v11.east) node [right,black,midway,text width=2.8cm,text height=1.2cm] {\textcolor{black}{Proglottids/eggs \\shed with stools}};
%
%\draw [-latex new, arrow head=2mm] [black,thick] (v8.west) -- (v6.south);
%\draw [-latex new, arrow head=2mm] [black,thick] (v11.west) -- (v12.east);
%\draw [-latex new, arrow head=2mm] [black,thick] (v6.north) -- (v2.west);
\end{tikzpicture}}
\caption{text}
\label{fig:spectralcomp}
\end{figure}


\newpage{\thispagestyle{empty}\cleardoublepage}
\chapter{General Discussion}\label{chap:gendiscussion}
\lettergroup{\thechapter}

text

\newpage{\thispagestyle{empty}\cleardoublepage}
\chapter{Bibliography}
\lettergroup{\thechapter}
\begingroup 
\def\chapter#1{} % needed because the thumb marks were not triggered for the first page of the biblio...
\small
\singlespace
\bibliographystyle{elsarticle-harv_M2}%elsarticle-harv_M2
\bibliography{ref} 
\endgroup
\onehalfspace

\newpage{\thispagestyle{empty}\cleardoublepage}
\chapter{Summary} 
\lettergroup{\thechapter}

\newpage{\thispagestyle{empty}\cleardoublepage}
\chapter{Samenvatting} 
\lettergroup{\thechapter}


\newpage{\thispagestyle{empty}\cleardoublepage} 

\appendix 
\chapter{Appendix}
\lettergroup{\thechapter}


\end{document}
